\documentclass{article}
\usepackage{/globals/boxes}
\usepackage{/globals/math}
\begin{document}
\section{Gleichungen}
Werden zwei Terme gleich gestellt, so wird eine Gleichung aufgestellt.

Sie zu lösen ist es ihre Lösungsmenge zu finden; die Menge an Werten der Variablen, für welche beide Seiten der Gleichung den selben Wert annehmen.

\subsection{Die Lösungsmenge}
Die Lösungsmenge \tL ist eine Menge, dessen Elemente die Lösungen sind. Gibt es keine Lösung, kann die Menge auch leer sein.
\[
 \mL = \{x\} \dtext{oder}
 \mL = \{x;y\} \dtext{oder}
 \mL = \{\}
\]

\subsection{Äquivalenzumformungen}
Einfache Gleichungen können gelöst werden, indem beide Terme äquivalent zueinander umgeformt werden; entweder äquivalent zu sich selbst, durch einfaches Umstellen mit bekannten Regeln, oder indem auf beiden Seiten die gleiche Rechnung durchgeführt wird.

Dabei muss darauf geachtet werden, dass bei keinen Schritt Information verloren geht, zum Beispiel durch eine Multiplikation mit Null, eine Division durch Null, oder der annahme dass ${1^2 = (-1)^2 \Leftrightarrow 1 = -1}$ oder ähnliches.

\subsection{p-q-Formel}
Ist eine quadratische Gleichung ohne Faktor vor dem ersten Glied gegeben, so kann die p-q-Formel angewand werden.
\[
 x^2 + px + q = 0
 \Rightarrow
 x_{1;2} = -\frac{p}{2} \pm \sqrt{\frac{p^2}{4}-q}
\]
Die Diskrimante $\frac{p^2}{4}-q$ kann einem hier verranten wie viele Nullstellen es gibt; so viele, wie dessen Wurzel Ergebnisse hat.

Oftmals liegt die gegebene Gleichung nicht in der obigen Form vor. Dann kann sie durch Äquivalenzumformungen, durch das Ausklammern und den Satz vom Nullprodukt, durch Substitution, oder durch Polynomdivision in diese gebracht werden.

Bei der Substitution können sich neue Variablen ausgedacht werden, dessen Werte Teile der Gleichung sind, so gewählt, dass mit den neuen Variablen die Gleichung in der obigen Form vorliegt.

\subsection{Als Produkt darstellen}
Sind die Nullstellen einer quadratischen Gleichung bekannt, so kann diese als Produkt dargestellt werden.
\[
 \prod_{k=0}^n x-x_k
\]
\end{document}